% 光的干涉（高中）
% keys 波的叠加
% license Usr
% type Tutor

\subsection{双缝干涉}
\subsubsection{理论推导}
\subsubsection{实验装置}
一个经典的双缝干涉实验装置如下图所示。光先通过单缝，成为柱面波，我们可以画出一系列同心圆，每一个圆代表一个相位，最前面的波面是光源的初始相位。也就是说，柱面波到达双缝的必是光源的初始相位。由此，我们在双缝处得到两个相位差为0的光源。
\begin{figure}[ht]
\centering
\includegraphics[width=12cm]{./figures/17f92705aa639da5.png}
\caption{} \label{fig_interf_1}
\end{figure}

\subsection{薄膜干涉}